\section{Nghiên cứu liên quan và mốc tham chiếu}
\label{sec:related_work}

\subsection{Mô hình phân giai đoạn giấc ngủ từ EEG}

Các phương pháp học sâu cho phân giai đoạn giấc ngủ thường kết hợp bộ trích đặc trưng từ tín hiệu với một mô hình mô tả quan hệ theo thời gian. DeepSleepNet sử dụng CNN và BiLSTM để đồng thời học hình thái của epoch và ngữ cảnh chuỗi~\cite{supratak2017deepsleepnet}. Các hướng tiếp cận sau đó khai thác cơ chế chú ý hoặc self-attention để tăng khả năng biểu diễn quan hệ dài hạn~\cite{eldele2021attnsleep,phan2022sleeptransformer}.

SleepInceptionNet khảo sát có hệ thống tín hiệu thô, phổ và biểu diễn thời gian--tần số, sau đó chọn mô hình Inception dùng biến đổi wavelet liên tục trên EEG trung tâm đơn kênh để chấm từng epoch~\cite{haghayegh2023sleepinceptionnet}. Kết quả đó củng cố vai trò của EEG đơn kênh, đồng thời cho thấy biểu diễn và tiền xử lý đầu vào là một phần của định nghĩa pipeline chứ không phải chi tiết trung tính.

Các chỉ số giữa những công trình khác nhau không thể xếp hạng trực tiếp nếu khác bộ dữ liệu, kênh EEG, quy tắc tạo cửa sổ hoặc cách chia đối tượng. Vì vậy, báo cáo chỉ thực hiện so sánh nội bộ giữa các cấu hình trong cùng một giao thức, thay vì suy ra thứ hạng giữa các bài báo.

Các nghiên cứu chuyển kịch bản đã cho thấy transfer learning có thể cải thiện phân giai đoạn giấc ngủ EEG khi miền nguồn và đích khác nhau~\cite{he2023crossscenario}. ADAST dùng attention riêng theo miền, căn chỉnh đối kháng và self-training bằng pseudo-label từ dữ liệu đích không nhãn~\cite{eldele2023adast}; mô hình được cập nhật bằng miền đích. E6 trong báo cáo hiện tại chỉ ước lượng thống kê chuẩn hóa ở cấp bản ghi và không cập nhật trọng số, nên không phải đối chứng trực tiếp cho phương pháp thích nghi đã học của ADAST. Các phương pháp máy học truyền thống cũng có thể cạnh tranh khi đặc trưng và giao thức được xây dựng cẩn thận~\cite{vdd2023traditional}. Những công trình này đặt lựa chọn pipeline, không chỉ độ sâu mạng, vào trung tâm của đánh giá thực nghiệm.

Davidson và cộng sự khảo sát lại tiêu chí chấm N3 trên 2.913 bản ghi SHHS, nhấn mạnh mối liên hệ giữa ngưỡng sóng chậm, tuổi và giới trong kết quả chấm~\cite{davidson2025n3}. Đây là bối cảnh quan trọng khi diễn giải lỗi N3 ngoài miền, nhưng nghiên cứu về tiêu chí chấm không chứng minh nguyên nhân của lỗi chuyển Sleep-EDF sang SHHS1 trong báo cáo này.

\subsection{Mốc tham chiếu và hướng thay thế}

ZleepAnlystNet là mốc tham chiếu quan trọng của nghiên cứu, với quy trình gồm bộ trích đặc trưng CNN và mô hình chuỗi BiLSTM~\cite{jadhav2024zleepanlystnet}. Ý tưởng thay thế được khảo sát theo hai hướng. Thứ nhất, mạng tích chập theo thời gian có khả năng xử lý chuỗi song song và mô hình hóa quan hệ theo trường tiếp nhận mở rộng~\cite{bai2018tcn}. Thứ hai, ResNet-1D sử dụng kết nối dư để xây dựng biểu diễn tín hiệu sâu hơn~\cite{he2016resnet}.

Hai hướng thay thế không mặc nhiên bảo đảm hiệu năng cao hơn. TCN có thể đem lại lợi ích vận hành nhưng vẫn cần được kiểm tra trên cùng đối tượng; embedding ResNet có thể hữu ích cho mô hình chuỗi mà không nhất thiết tạo ra các cụm hình học tách biệt. E1--E0 giữ encoder/đặc trưng nhưng thay cả cấu hình mô hình chuỗi và lịch huấn luyện. E2--E1 thay encoder, gói ngữ cảnh C/P/N, số chiều đặc trưng và cách huấn luyện/chọn encoder. Cả hai vì vậy được diễn giải ở cấp cấu hình, không phải tác động kiến trúc riêng lẻ.

\subsection{Khoảng trống được kiểm tra}

Nghiên cứu kết hợp ba góc nhìn thường được trình bày riêng: hiệu năng bắt cặp của các cấu hình pipeline, chi phí vận hành và lỗi theo lớp khi chuyển cohort. Phép chia theo đối tượng và thời điểm khóa test được giữ chung; các khác biệt về huấn luyện và chọn encoder được khai báo để người đọc biết chính xác giới hạn của từng đối chiếu.

Bên cạnh hiệu năng trong miền, nghiên cứu xem xét liệu thứ hạng quan sát trên Sleep-EDF và các cải thiện tổng thể còn đi cùng chất lượng nhận diện từng lớp trên SHHS1 hay không. Báo cáo không đánh giá một thuật toán domain adaptation mới; kết quả E6 và phân tích kênh lỗi cung cấp các giả thuyết có mục tiêu cho nghiên cứu tiếp theo, không xác nhận rằng một chiến lược thích nghi hoặc thu nhãn cụ thể là tối ưu.

\subsection{Nguyên tắc diễn giải}

Không bác bỏ giả thuyết không không đồng nghĩa với chứng minh tương đương. Tương tự, một chỉ số hình học hoặc một phương pháp ước lượng tầm quan trọng không thể tự nó xác định nguyên nhân của chênh lệch dự đoán. Báo cáo vì vậy ưu tiên chênh lệch hiệu năng, khoảng tin cậy, kiểm định bắt cặp và phạm vi áp dụng của từng kết quả.
